\documentclass[aspectratio=169, 10pt]{beamer}

% ======================================================================
% PAQUETES Y TEMA
% ======================================================================
\usepackage[utf8]{inputenc}
\usepackage[spanish,es-tabla]{babel}
\usepackage{graphicx}
\usepackage{xcolor}
\usepackage{booktabs}
\usepackage{siunitx}
\usepackage[american]{circuitikz}
\usepackage{listings}

\usetheme{Boadilla}
\usecolortheme{seahorse}

\sisetup{output-decimal-marker = {,}, per-mode = symbol}

\lstset{
    language=Python,
    basicstyle=\ttfamily\scriptsize,
    keywordstyle=\color{blue}\bfseries,
    commentstyle=\color{gray}\itshape,
    backgroundcolor=\color{gray!5},
    frame=single,
    breaklines=true
}

% ======================================================================
% METADATOS
% ======================================================================
\title[Filtro Twin-T Notch]{Semana 3 - Sesión 2\\Análisis Fasorial: Filtro Twin-T Notch}
\subtitle{Circuitos Electrónicos II (IMT-221)}
\author[M.Sc. Ing. B. Quiroga]{M.Sc. Ing. Bernardo Quiroga Turdera}
\institute[UCB Tarija]{Universidad Católica Boliviana ``San Pablo''}
\date{Semestre 2-2026}

\begin{document}

\begin{frame}
    \titlepage
\end{frame}

\begin{frame}{Enfoque Pedagógico: Análisis Fasorial Riguroso}
    \textbf{¿Por qué no usamos la variable $s$ de Laplace?}
    \vspace{0.3cm}
    \begin{itemize}
        \setlength{\itemsep}{5pt}
        \item \textbf{Intuición Física:} Evitamos abstracciones algebraicas prematuras.
        \item \textbf{Conexión Real:} Trabajamos con $\omega = 2\pi f$. Cada término se traduce directamente a perillas del generador (Hz) y lecturas de osciloscopio (V, $\mu$s).
        \item \textbf{Rigor Matemático:} Cumplimos con ABET (Student Outcome 1) mediante impedancias complejas y Leyes de Kirchhoff exactas.
    \end{itemize}
\end{frame}

\begin{frame}{El Problema Industrial en el PFI}
    \begin{columns}[T] % Alineación superior para evitar desplazamientos
        \begin{column}{0.40\textwidth}
            \textbf{Situación:}
            \begin{itemize}
                \setlength{\itemsep}{3pt}
                \item Portadora útil: \qty{10}{\kilo\hertz}
                \item Ruido acoplado de red: \qty{50}{\hertz}
            \end{itemize}
            \vspace{0.3cm}
            \textbf{Solución:}\\
            Filtro supresor de banda (Notch) sintonizado a $f_0 = \qty{50}{\hertz}$.
        \end{column}
        
        \begin{column}{0.60\textwidth}
            % DIAGRAMA ESCALADO PARA ENCAJAR PERFECTAMENTE
            \begin{center}
                \begin{circuitikz}[scale=0.55, transform shape]
                    \draw (0,2) to[sV, l=$V_{in}$] (0,0) node[ground]{};
                    \draw (0,2) to[short, -*] (1.5,2);

                    \draw (1.5,2) -- (1.5,4.5);
                    \draw (1.5,2) -- (1.5,-0.5);

                    \draw (1.5,4.5) to[R, l=$R$] (4.5,4.5) coordinate(nA);
                    \draw (nA) to[R, l=$R$] (7.5,4.5);
                    \draw (nA) to[C, l=$2C$, *-] (4.5,2.5) node[ground]{};
                    \node[above=0.15cm] at (nA) {\textbf{Nodo A}};

                    \draw (1.5,-0.5) to[C, l=$C$] (4.5,-0.5) coordinate(nB);
                    \draw (nB) to[C, l=$C$] (7.5,-0.5);
                    \draw (nB) to[R, l=$R/2$, *-] (4.5,-2.5) node[ground]{};
                    \node[above=0.15cm] at (nB) {\textbf{Nodo B}};

                    \draw (7.5,4.5) -- (7.5,2);
                    \draw (7.5,-0.5) -- (7.5,2);
                    \draw (7.5,2) to[short, *-o] (8.5,2) node[right]{\large $V_{out}$};
                \end{circuitikz}
            \end{center}
        \end{column}
    \end{columns}
\end{frame}

\begin{frame}{Deducción Matemática (Leyes de Kirchhoff)}
    Aplicando la Ley de Corrientes en cada nodo ($Z_R=R, Z_C=\frac{1}{j\omega C}$):
    \vspace{0.2cm}
    \begin{block}{Nodo A (Rama Paso-Bajo)}
        $\hat{V}_A = \frac{\hat{V}_{in} + \hat{V}_{out}}{2(1 + j\omega RC)}$
    \end{block}
    \begin{block}{Nodo B (Rama Paso-Alto)}
        $\hat{V}_B = \frac{j\omega RC (\hat{V}_{in} + \hat{V}_{out})}{2(1 + j\omega RC)}$
    \end{block}
    \vspace{0.2cm}
    Sumando corrientes en el nodo de salida $\hat{V}_{out}$:
    $$H(j\omega) = \frac{\hat{V}_{out}}{\hat{V}_{in}} = \frac{1 - \left(\frac{\omega}{\omega_0}\right)^2}{1 - \left(\frac{\omega}{\omega_0}\right)^2 + j 4 \left(\frac{\omega}{\omega_0}\right)} \quad \text{donde } \omega_0 = \frac{1}{RC}$$
\end{frame}

\begin{frame}{Comportamiento Físico y Asintótico}
    $$H(j\omega) = \frac{1 - (\omega/\omega_0)^2}{1 - (\omega/\omega_0)^2 + j 4 (\omega/\omega_0)}$$
    \vspace{0.2cm}
    \begin{itemize}
        \setlength{\itemsep}{8pt}
        \item \textbf{Resonancia ($\omega = \omega_0 \implies \qty{50}{\hertz}$):} Numerador cero $\implies \vert{}H\vert{} = 0\text{ V}$.
        \begin{itemize}
            \item \textit{Explicación Física:} Interferencia destructiva total (desfases de $+90^\circ$ y $-90^\circ$ de ambas ramas se anulan).
        \end{itemize}
        \item \textbf{Frecuencia Cero ($\omega \to 0$):} $\vert{}H\vert{} = 1 \implies \qty{0}{\decibel}$.
        \item \textbf{Portadora Alta ($\omega \gg \omega_0 \implies \qty{10}{\kilo\hertz}$):} $\vert{}H\vert{} \approx 1 \implies \qty{0}{\decibel}$.
    \end{itemize}
\end{frame}

\begin{frame}{El Peligro del Efecto de Carga ($R_L$)}
    Si conectamos una resistencia de carga finita $R_L$:
    $$H_{cargado}(j\omega) = \frac{1 - \left(\frac{\omega}{\omega_0}\right)^2}{1 - \left(\frac{\omega}{\omega_0}\right)^2 + j 4 \mathbf{\left(1 + \frac{R}{R_L}\right)} \left(\frac{\omega}{\omega_0}\right)}$$
    
    \vspace{0.3cm}
    \textbf{Conclusión vital:} \\
    Si $R_L$ es pequeña (ej. \qty{10}{\kilo\ohm}), el término imaginario crece, degradando el factor de calidad ($Q$) y arruinando la profundidad de la muesca (baja de $\qty{-50}{\decibel}$ a $\qty{-21}{\decibel}$). \textit{¡Siempre requiere un Buffer Op-Amp a la salida!}
\end{frame}

\begin{frame}{Resolución Numérica (Diseño Real)}
    \textbf{Datos:} $R = \qty{14.3}{\kilo\ohm}, C = \qty{220}{\nano\farad}$. Interferencia \qty{1.2}{\volt} a \qty{50}{\hertz}. Portadora \qty{150}{\milli\volt} a \qty{10}{\kilo\hertz}. Atenuación lograda en protoboard: $\qty{-44}{\decibel}$.
    \vspace{0.3cm}
    \begin{enumerate}
        \setlength{\itemsep}{6pt}
        \item \textbf{Sintonía:} $f_0 = \frac{1}{2\pi \cdot 14300 \cdot 220\times 10^{-9}} = \mathbf{50.59\text{ Hz}}$
        \item \textbf{Ruido Residual a \qty{50}{\hertz}:} $\hat{V}_{ruido} = 1.2 \cdot 10^{-44/20} = \mathbf{7.57\text{ mV}_{pico}}$ (casi nulo)
        \item \textbf{Transmisión de Portadora:} $\hat{V}_{port} = 150 \times 0.9998 = \mathbf{149.97\text{ mV}_{pico}}$ (pérdida nula)
        \item \textbf{Carga de \qty{10}{\kilo\ohm}:} $\vert{}H\vert{} = \mathbf{-21.3\text{ dB}}$ (falla de diseño por carga)
    \end{enumerate}
\end{frame}

\begin{frame}{Tarea Pre-Laboratorio: ¡Emparejamiento Físico!}
    \textbf{Para el montaje del próximo lunes (15 minutos máximo):}
    \vspace{0.3cm}
    \begin{itemize}
        \setlength{\itemsep}{5pt}
        \item \textbf{Multímetro:} Medir y rotular $R_1, R_2 \approx \qty{14.3}{\kilo\ohm}$ (idénticos $< 0.5\%$).
        \item \textbf{Trimpot:} Ajustar $R_3 \approx \qty{7.15}{\kilo\ohm}$ (exactamente la mitad).
        \item \textbf{Capacímetro:} Emparejar $C_1, C_2 \approx \qty{220}{\nano\farad}$ y armar paralelo $C_3 \approx \qty{440}{\nano\farad}$.
        \item \textbf{Entregables Virtuales:} Archivo \texttt{.asc} de LTSpice con directivas de barrido y script \texttt{lab2\_notch\_phasor.py} funcional.
    \end{itemize}
\end{frame}

\end{document}